% !PRINTED 432
\weberpage{432}
%<<CONT
und aus
\weq{10}{p\,\alpha'\,\delta-\beta\,\gamma = 1}
folgt sodann
\weq{11}{y^2+m\,x^2 = 4.}

Der Wert $x=0$ ist auszuschließen, weil sonst $\beta=0$ sein müßte, was der Gleichung (10) widerspricht, und wir können unbeschadet der Allgemeinheit $x$ positiv annehmen. Da überdies $m$ nach dem Modul 4 entweder $\equiv 0$ oder $\equiv 3$ sein muß, so bleiben zur Erfüllung von (11) nur folgende zwei Möglichkeiten übrig:
\[
1.\quad m=3,\quad x=1,\quad y=\pm 1.
\]

Da es für die Diskriminante $-3$ nur eine Formenklasse gibt, so können wir $A=1,\,B=1,\,C=1$ annehmen, d.\,h. für $\omega$ die imaginäre dritte Einheitswurzel $e^{\frac{2\pi i}{3}}$ setzen, und erhalten unter den Größen (3) drei äquivalente, indem wir $y=+1$ und $=-1$ annehmen: % [sic] Im Druck fehlt vor dem zweiten Gleichheitszeichen die Variable y.
\weq{12}{p\,\omega,\quad \frac{\omega}{p},\quad \frac{\omega+1}{p},}
wie auch aus den Gleichungen
\[
\frac{\omega}{p} = \frac{-1}{p+p\,\omega}, \quad \frac{\omega+1}{p} = \frac{-1}{p\,\omega}
\]
erkannt wird, aus denen die Äquivalenz der drei Größen (12) evident ist.
\[
2.\quad m=4,\quad x=1,\quad y=0.
\]

Weil es auch hier nur eine Formenklasse gibt, so können wir $B=0,\,A=C=1$, d.\,h. $\omega=i$ annehmen, und finden die zwei äquivalenten Werte:
\[
p\,\omega,\quad \frac{\omega}{p},
\]
wie auch aus der Gleichung
\[
p\,\omega = \frac{-p}{\omega}
\]
folgt.

Es seien ferner zwei der Werte (3)
\[
\frac{\omega+c}{p}, \quad \frac{\omega+c'}{p}
\]
äquivalent, so daß eine Gleichung besteht:
\weq{13}{\frac{\omega+c}{p} = \frac{\gamma\,p+\delta\,(\omega+c')}{\alpha\,p+\beta\,(\omega+c')}, \quad \alpha\,\delta-\beta\,\gamma = 1,}
%CONT>>
